长度为 \(L>0\)、线密度为 \(\mu>0\)、张力为 \(T>0\) 的理想细弦两端固定。弦在正中心 \(x=L/2\) 与质量 \(M>0\) 的微小质点固定连接，质点只作横向振动。弦的两侧仍满足 \(y_{tt}=c^2 y_{xx}\)，其中 \(c=\sqrt{T/\mu}\)，质点处位移连续但斜率可有跳跃。忽略阻尼。
(1) 由质点牛顿第二定律导出中心的斜率跳跃边界条件；(2) 利用中心反射对称，将本征模分为中心位移为零的反对称模和中心位移不为零的对称模，分别推导特征频率方程；(3) 证明反对称模频率不受 \(M\) 影响，并用动能与势能表达式说明所有非零本征频率都为正；(4) 求 \(M\to0^+\) 和 \(M\to\infty\) 时最低对称模的频率极限，解释大质量极限出现的低频软模。
